% !PRINTED none
\begin{center}
\gesperrt{Nineteenth Chapter.}

\gesperrt{Relations of a Field to its Divisors.}

\rule{0.2\linewidth}{0.4pt}

\S.\,159.

\gesperrt{Relative Norms.}
\end{center}

In this chapter we shall consider the modifications that occur in the definitions and theorems of the theory of algebraic numbers when, in place of the \gesperrt{absolute domain of rationality}, that is, of the field of the rational numbers, an arbitrary algebraic number field is substituted. In this connection some important supplements will also result for the general theory of prime factors, which find their expression in a series of theorems established by \gesperrt{Dedekind} and \gesperrt{Hilbert}${}^{1}$).

Let $R$, then, be an algebraic number field of degree $m$, and let
\weq{1}{f(\Theta) = \Theta^n + \alpha_1\Theta^{n-1} + \cdots + \alpha_n = 0}
be an equation of degree $n$, rational and irreducible in $R$. The field $\Omega = R(\Theta)$ is an algebraic field \gesperrt{over} $R$, which with respect to $R$ is of degree $n$ (Bd. I, \S.\,142). In the absolute domain of rationality $\Omega$ is of degree $mn$, and if $\varrho$ is a primitive number of the field $R$, then
\weq{2}{\Theta^s \varrho^t, \qquad \begin{aligned} s &= 0, 1, \ldots, n-1\\ t &= 0, 1, \ldots, m-1 \end{aligned}}
is a basis of the field $\Omega$; for, owing to the irreducibility of the
%CONT>>

${}^{1}$) \gesperrt{Dedekind}, ``Ueber die Discriminanten endlicher K\"orper'', Abhandlungen der G\"ottinger Gesellschaft der Wissenschaften (1882). ``Zur Theorie der Ideale'', Nachrichten der Gesellschaft der Wissenschaften zu G\"ottingen, 1894, Nr. 4. \gesperrt{Hilbert}, ``Grundz\"uge einer Theorie des \gesperrt{Galois'schen} Zahlk\"orpers''. Ibid. 1894, Nr. 3.
